\section{Introduction}\label{sec:intro}

Zero-knowledge (ZK) proofs now guard real money on public blockchains.
Privacy pools release deposits to whoever proves membership, rollups settle
state on a proof of correct execution, identity systems mint credentials to
the holder of a proved attribute, and zkVM markets pay provers for proved
computations. In each case a smart contract, the \emph{consumer}, accepts a
proof and then acts. Much work makes sure that proofs are sound, circuits
fully constrained and verifiers correctly
generated~\cite{sok-snark-bugs,picus,circomspect,zkap}. Much less asks a
simpler question: does the contract's action follow from what was proved?

Often it does not. A proof certifies a \emph{statement}: the public inputs
and what the circuit fixes about them. The consumer performs an
\emph{effect}, which depends on a recipient, an amount, a state root, a
replay key or a verification key. When an effect value is not fixed by the
statement, whoever submits the proof chooses it, and since proofs are public
once broadcast, that can be an attacker who copies a victim's proof and
changes one argument. We call this a \emph{binding flaw}: the proof is valid,
the verifier is correct, and the effect is still unauthorized.
It lives at the \emph{seam} between the proof side (circuit or zkVM guest) and
the consumer. Figure~\ref{fig:motivating} shows one from an audited proof
marketplace.

Binding flaws cause real losses. In June 2026 an attacker drained \$2.28M
from the Aztec Connect rollup contract without breaking any cryptography:
every batch carried a valid proof, but the contract settled as many rows as
a calldata field declared, and the proof did not bind that
field~\cite{aztec-connect-hack}. ZKSwap's bridge lost \$5M in 2025 to an
exit verifier that returned \emph{true} before checking
anything~\cite{zkswap-hack}. In a public record of ZK incidents from 2023 to
2026, every circuit-soundness bug in an on-chain application was caught
before exploitation, while each quantified loss we could trace sits at the
verifier or at the contract that consumes it~\cite{zkp-hacks-record}.

\begin{figure}[t]
\begin{lstlisting}[language=Java]
function fulfill(Fulfillment calldata fill, AssessorReceipt calldata r) public {
  verifyDelivery(fill, r);  // assessor journal commits fill.requestDigest
  ...
}
function _fulfillAndPay(Fulfillment calldata fill, address prover) internal {
  RequestId id = fill.id;                       // free calldata field
  (address client, uint32 idx) = id.clientAndIndex();
  ... _fulfillAndPayNeverLocked(id, client, idx, fill.requestDigest, ...);
}  // charges client of fill.id; nothing checks fill.id == id(requestDigest)
\end{lstlisting}
\caption{A binding flaw in an audited zkVM proof marketplace (abridged
from the vulnerable commit; since fixed). The proof binds the request \emph{digest}; the contract charges the
request named by a calldata \emph{identifier}. An attacker proves its own
cheap request but names a victim's, which is charged and marked fulfilled.
Each side is correct on its own; the flaw is the missing equality between
them.}
\label{fig:motivating}
\end{figure}

\para{Common, severe, and invisible to today's tools}
We collected 8,578 public audit reports and labeled 264 proof-related
findings. 87 are core binding flaws, spread over 43 of 68 project families;
78\% of the rated ones are Critical or High, and 55 of the 311 in-scope ZK
application audits report at least one (Section~\ref{sec:study}). Yet
Slither and Mythril report none of the audited flaws of our test set at the
right location and mechanism, circomspect and Picus cannot analyze most
targets because most binding flaws are not in a circuit, and a fuzzer with
generic properties finds none (Section~\ref{sec:eval-rq2}). The reason is
structural. A binding flaw breaks no rule of the contract and no constraint
of the circuit; it breaks an intended correspondence between the two that
neither artifact states. This is also why it is not input validation, which
checks a value against the program's own rules: given the missing
correspondence as a hand-written invariant, the same fuzzer finds 5 of the
13 flaws it could run on, against 0 without it.

\para{What an LLM agent gets wrong}
LLMs read Solidity, Circom, Noir and Rust alike, so an agent can in
principle follow a value from a guest's journal to a contract's storage
write. We build \tool{}, an agent that must fill both ledgers row by row
under coverage checks, and the ledgers let us ask a question that is usually
unanswerable for LLM tools: when it misses a flaw, had it never seen it, or
had it seen it and lost it? The answer surprised us. On 28 development flaws the agent
misses 16; we recovered six, and none of the six had been unseen. Three were reported on the \emph{other side} of the
proof boundary, at the contract when the root cause lay in the guest or the
reverse; three sat in the agent's ledgers as \emph{unbound} and were
withdrawn for reasons the threat model does not support, such as ``the
consuming contract is not in this repository''. Undoing both with one
tool-less call is far more effective, and far more specific, than teaching
the agent more: a pitfall catalogue, a second tool-using look at every
verdict, input inventories and targeted probes each failed a specificity
rule fixed in advance, and those that raised recall did so only by adding
false findings.

\para{Contributions}
\begin{itemize}
\item \textbf{A definition and taxonomy} of binding flaws: two ledgers under
an intended policy, three conditions, and five patterns with
signature-misuse analogues (Section~\ref{sec:definition}).
\item \textbf{To our knowledge, the first measurement study} of the class:
8,578 reports, 264 labeled findings and 87 core binding flaws, with
prevalence, severity, patterns, trends and fixes (Section~\ref{sec:study}).
\item \textbf{\tool{}}, a ledger-guided LLM agent with a strictly additive
post-stage that re-adjudicates discarded suspicions and reports each flaw at
every root-cause location, and an execution loop that attaches attack tests
to warnings (Section~\ref{sec:design}).
\item \textbf{An evaluation} on audited flaws, patched versions, hard
negatives and injected mutants against static, circuit, fuzzing and LLM
baselines; an anatomy of agent misses with four negative results under a
fixed keep rule; \DEMO{and a held-out evaluation and three new flaws
confirmed by executed tests} (Section~\ref{sec:eval}). We release the
dataset, benchmarks and tool.
\end{itemize}
